\documentclass[11pt,a4paper]{article}
\usepackage{fontspec,amsmath,amssymb,amsthm,microtype}
\setmainfont{lmroman10-regular.otf}[BoldFont=lmroman10-bold.otf,ItalicFont=lmroman10-italic.otf,BoldItalicFont=lmroman10-bolditalic.otf]
\setmonofont{lmmono10-regular.otf}[Scale=MatchLowercase]
\usepackage[margin=25mm]{geometry}
\usepackage[unicode,hidelinks]{hyperref}
\setlength{\emergencystretch}{3em}
\setlength{\parskip}{0.4em}
\setlength{\parindent}{0pt}
\allowdisplaybreaks[2]
\raggedbottom
\title{\texorpdfstring{The actual canonical KLMC coupling on each initial transport plan}{The actual canonical KLMC coupling on each initial transport plan}}
\author{LangevinLean mathematical exposition}
\date{}
\begin{document}\maketitle
\noindent Source module: \texttt{Kinetic\allowbreak{}Langevin/\allowbreak{}Canonical\allowbreak{}Coupled\allowbreak{}Segments.\allowbreak{}lean}. The accompanying HTML pairs this mathematical exposition with the exact Lean statements and records the scope of the formalization.


Fix any finite-cost coupling of the initial position law and the target position law. The probability space is its explicit product with an independent standard Gaussian velocity and the canonical vector Brownian source. This module constructs the numerical chain, the exact reference segments and the auxiliary exact segments on that same space, and proves every analytic segment requirement used by the quantitative KLMC theorem.


\section{\texorpdfstring{Specify the concrete probability space and source projection}{Specify the concrete probability space and source projection}}


The phase variable is velocity first. The sample space contains the two coupled initial positions, one shared Gaussian velocity and one vector Brownian source. The measure is the actual initial-coupling product extension; probability follows from its first marginal and from the Gaussian/\allowbreak{}Brownian probability factors. The Brownian-source projection preserves its measure. Pullback along that projection gives a continuous linear map on L\ensuremath{{}^{2}}, and its representative theorem identifies it almost everywhere with evaluation on the actual product sample.



{\raggedright\noindent Lean declarations: \texttt{H}, \texttt{Sample}, \texttt{probability}, \texttt{probability\_\allowbreak{}is\allowbreak{}Probability\allowbreak{}Measure}, \texttt{coupling\allowbreak{}Probability}, \texttt{brownian\allowbreak{}Lift}, \texttt{brownian\allowbreak{}Lift\_\allowbreak{}coe}.\par}

\section{\texorpdfstring{Construct the common Brownian input and filtered innovation}{Construct the common Brownian input and filtered innovation}}


The scaled source lift includes the physical amplitude \(\sqrt{2\gamma}\). At iteration \(n\), the input is the actual vector Brownian increment on \((nh,nh+t]\), and the filtered input is the exponential response to that same increment. The unscaled innovation pair consists of the exact velocity and position Wiener responses at the step endpoint. Both input paths are continuous in L\ensuremath{{}^{2}}, and the additive input starts at zero. These follow from the source lift's continuity and linearity and the proved vector-response properties.



{\raggedright\noindent Lean declarations: \texttt{input\allowbreak{}Lift}, \texttt{input}, \texttt{filtered}, \texttt{innovation}, \texttt{continuous\_\allowbreak{}input}, \texttt{continuous\_\allowbreak{}filtered}, \texttt{input\_\allowbreak{}zero}.\par}

\section{\texorpdfstring{Identify the exact noise equations and implemented numerical step}{Identify the exact noise equations and implemented numerical step}}


The constructed filtered input satisfies \(B_n(t)=W_n(t)-\gamma\int_0^tB_n(s)ds\). Its endpoint is the scaled velocity innovation, and its actual time primitive is the scaled positional innovation. Lean transports the vector Brownian equations through the continuous linear source lift and commutes that lift with Bochner integration. Induction on the actual recursions then proves that the exponential-input chain is exactly \texttt{klmc\allowbreak{}Chain}, with the canonical correlated innovation pair and the prescribed fluctuation-dissipation scaling.



{\raggedright\noindent Lean declarations: \texttt{filtered\_\allowbreak{}integral\_\allowbreak{}equation}, \texttt{integral\_\allowbreak{}filtered}, \texttt{filtered\_\allowbreak{}endpoint}, \texttt{exponential\allowbreak{}Input\allowbreak{}Chain\_\allowbreak{}eq\_\allowbreak{}klmc\allowbreak{}Chain}.\par}

\section{\texorpdfstring{Connect Brownian L\ensuremath{{}^{2}} innovations with the generic output-\allowbreak{}law construction}{Connect Brownian L² innovations with the generic output-law construction}}


The L\ensuremath{{}^{2}} innovation pair has the same representatives as the pointwise canonical Brownian innovation. The proof pulls back the previously established vector-response representative theorem to the coupling product and combines the two component identities. Comparing with the generic product-law layer's representatives gives equality of the actual L\ensuremath{{}^{2}} pairs. Thus the numerical chain analyzed here has the coupling-independent output law already proved for the implemented KLMC chain.



{\raggedright\noindent Lean declarations: \texttt{innovation\_\allowbreak{}coe}, \texttt{innovation\_\allowbreak{}eq\_\allowbreak{}extension\allowbreak{}Innovation}.\par}

\section{\texorpdfstring{Build the actual numerical and stationary-\allowbreak{}reference initial states}{Build the actual numerical and stationary-reference initial states}}


The two initial phases share exactly one Gaussian velocity variable. The numerical phase uses the first coupled position; the reference phase uses the second. The reference representative is proved to be the explicit pair \((V,X_\mathrm{ref})\). Product independence and the two marginal laws prove that its actual law is \(\mathcal N(0,I_d)\otimes\pi\). The finite target second moment and finite coupling cost supply the required L\ensuremath{{}^{2}} initial positions.



{\raggedright\noindent Lean declarations: \texttt{initial\allowbreak{}Velocity}, \texttt{numerical\allowbreak{}Initial}, \texttt{reference\allowbreak{}Initial}, \texttt{reference\allowbreak{}Initial\_\allowbreak{}coe}, \texttt{reference\allowbreak{}Initial\_\allowbreak{}law}.\par}

\section{\texorpdfstring{Construct all three chains and segment families}{Construct all three chains and segment families}}


The numerical chain applies the actual lifted original gradient and the constructed correlated innovations. The reference chain stitches exact selected kinetic segments, each driven by the actual canonical Brownian increment for that iteration. The corresponding reference segments start from these actual reference states. The auxiliary exact segment starts instead from the actual numerical state and receives the identical additive input. The drift is the genuine pointwise gradient lifted to L\ensuremath{{}^{2}}; no substitute potential on the L\ensuremath{{}^{2}} space is assumed.



{\raggedright\noindent Lean declarations: \texttt{numerical\allowbreak{}Chain}, \texttt{reference\allowbreak{}Chain}, \texttt{reference\allowbreak{}Segment}, \texttt{auxiliary\allowbreak{}Segment}.\par}

\section{\texorpdfstring{Prove every analytic premise of the KLMC coupling theorem}{Prove every analytic premise of the KLMC coupling theorem}}


The bundled specification supplies continuity of auxiliary paths, reference paths and filtered inputs; exact initial states; reference endpoint matching; the implemented exponential update; both actual common-input integral equations; the exponential mild identity; the kinematic primitive identity; and equal initial velocities. Each field follows from the concrete definitions and earlier proved equations. The mild and kinematic fields are derived internally. The only remaining probabilistic input for a final stationary-error theorem is the stationary law of the actual global reference diffusion, proved separately and connected to these stitched segments by uniqueness.



{\raggedright\noindent Lean declarations: \texttt{segments\_\allowbreak{}spec}.\par}

\end{document}
